\section{Population maximum likelihood and a repaired uniqueness theorem}
\label{sec:population}

The finite failures of Section~\ref{sec:finite} have two different causes, and
separating them is what motivates the model of this section.

\emph{Structural non-identifiability} (Example~\ref{ex:ABAB}) is the failure to
determine the genome from the read distribution at all. Bridging conditions
address it, and Theorem~\ref{thm:rigidity} shows how far they go on the
same-length slice.

\emph{Finite-sampling fluctuation} is a different phenomenon. Even a candidate
that is structurally admissible can beat the truth because the realized sample
is unrepresentative. The next example isolates this mechanism in its purest
form.

\begin{example}[Finite frequencies can beat structural admissibility]
\label{ex:frequencies}
Let $S=AABBC$ ($G=5$) be a primitive circular genome and $L=3$. Realize the
starts $(0,3)$, giving the observed reads $\{AAB,BCA\}$. Let $D=AABC$ and note
$|D|=4$. Both $S$ and $D$ are primitive and satisfy P1 (no repeated
$(L-1)=2$-mer). The realized sample satisfies the truth-side
information-feasibility condition, yet under the candidate-intrinsic exact
multinomial objective
\[
  \frac{\Lik_{\mathrm{ex}}(D;x)}{\Lik_{\mathrm{ex}}(S;x)}
  = \frac{(1/4)^2}{(1/5)^2} = \frac{25}{16} > 1 .
\]
The wrong genome wins because the realized sample consists only of reads to
which the shorter candidate assigns probability $1/4$ while the truth assigns
$1/5$. Structural admissibility has not failed; the empirical frequencies have
fluctuated.
\statusline{\computed{}, exact rational arithmetic. This is the project's
finite witness for the sampling mechanism; it is not a counterexample to any
historical statement, since the candidate restrictions are project-level.}
\end{example}

The lesson is that strengthening the repeat axioms again would target the wrong
mechanism. The repair that removes empirical fluctuation is to replace the
finite sample by its population limit.

\subsection{Population likelihood}

\begin{definition}[Population read distribution and likelihood]
\label{def:population}
For a circular word $D$ and read length $L\le|D|$, the population read
distribution is
\[
  p_D(w) \;=\; \frac{\spec_L(D)(w)}{|D|},
\]
a probability distribution on length-$L$ words. For a fixed truth $S$, the
\emph{population log-likelihood} of a candidate $D$ is
\[
  \ellpop_S(D) \;=\; \sum_{w} p_S(w)\log p_D(w),
\]
with $\ellpop_S(D)=-\infty$ if some word with $p_S(w)>0$ has $p_D(w)=0$.
\statusline{\projchoice. This is the infinite-read (population) analogue of the
Medvedev--Brudno exact objective; it is a project-level formulation motivated by
the finite analysis, not an assumption of either historical paper.}
\end{definition}

\begin{lemma}[Gibbs/KL identity; the truth is always a population maximizer]
\label{lem:gibbs}
For any candidate $D$,
\[
  \ellpop_S(D)-\ellpop_S(S) = -\KL(p_S\,\|\,p_D) \;\le\; 0,
\]
with equality if and only if $p_D=p_S$. Consequently, over \emph{any} candidate
class containing $S$, the truth is a population maximum-likelihood genome, with
no repeat condition required.
\end{lemma}

\begin{proof}
$\ellpop_S(S)=\sum_w p_S(w)\log p_S(w)$, so
$\ellpop_S(D)-\ellpop_S(S)=\sum_w p_S(w)\log(p_D(w)/p_S(w))=-\KL(p_S\|p_D)$,
where the right side is $+\infty$ precisely when $\ellpop_S(D)=-\infty$. Gibbs'
inequality \cite{cover2006} gives $\KL\ge0$ with equality iff the distributions
agree. The last claim is immediate since $S$ is in the class.
\end{proof}

Lemma~\ref{lem:gibbs} is the conceptual heart of the repair. At the population
level the maximizer statement is free; the entire content of a uniqueness theorem
is then \emph{spectrum identifiability}: for which candidate classes is
$p_D=p_S$ equivalent to $D\approx S$?

\subsection{P2 plus primitivity remove spectrum scaling}

Population equality compares \emph{normalized} spectra. The project uses two
candidate-intrinsic admissibility predicates to cross the gap to ordinary
spectrum equality.

\begin{definition}[P1 and P2]
\label{def:P1P2}
Fix the read length $L$. An oriented circular word $D$ satisfies
\begin{itemize}[leftmargin=*]
  \item \textbf{P1} if no length-$(L-1)$ word occurs more than once in $D$;
  \item \textbf{P2} if $D$ has no triple repeat of length $\ge L-1$ and every
        interleaved maximal repeat pair has a constituent of length
        $\le L-2$.
\end{itemize}
P1 is stronger than P2 (no repeated $(L-1)$-mer implies the P2 repeat bounds).
Both are project-level, candidate-intrinsic conditions; neither is a historical
assumption.
\statusline{\projchoice. The threshold in P2 is chosen to match the
complete-spectrum uniqueness theorem at $K=L-1$ below.}
\end{definition}

\begin{lemma}[Primitive P2 spectra have gcd one]
\label{lem:scaling}
Let $S$ be a primitive P2-admissible circular word and $L\ge2$ with
$L\le|S|$. Let $c=\spec_L(S)$ viewed as an integer vector and let
$g=\gcd\{c(w):c(w)>0\}$. Then $g=1$.
\end{lemma}

\begin{proof}[Proof sketch]
Suppose $g>1$. Work in the order-$(L-1)$ de Bruijn graph: vertices are
$(L-1)$-mers and each occurrence of a length-$L$ word is a directed edge
with multiplicity $c(w)$. Because $S$ spells a closed trail using every
edge exactly once, this edge multiset is balanced (in-degree equals
out-degree at every vertex) and weakly connected on its support.
Dividing every multiplicity by $g$ gives $c'=c/g$, still integer-valued.
Balance is a homogeneous linear condition, so it is preserved by division;
the support is unchanged ($c'(w)>0\iff c(w)>0$), so weak connectivity on
the support is preserved as well. Hence $c'$ is a balanced connected edge
multiset and spells an Eulerian closed trail $W$ with
$\spec_L(W)=c'$ and $|W|=|S|/g$.
Repeating $W$ gives $W^g$ with $|W^g|=|S|$ and
$\spec_L(W^g)=g\cdot\spec_L(W)=c=\spec_L(S)$: each cyclic length-$L$
window of $W$ lifts to exactly $g$ windows of $W^g$.
Thus $S$ and $W^g$ share the same ordinary $L$-spectrum.
$S$ satisfies P2, hence Ukkonen's condition at $K=L-1$, so
Theorem~\ref{thm:BBT} applied to $S$ forces $S\sim W^g$ up to rotation.
But $W^g$ is a nontrivial whole-genome power while $S$ is primitive, and
primitivity is rotation-invariant --- a contradiction. So $g=1$.
Crucially, uniqueness is applied only to $S$; nothing requires $W^g$
itself to satisfy P2.
\end{proof}

\begin{corollary}[Normalized equality implies ordinary equality for primitive P2]
Let $S,D$ be primitive P2-admissible with $p_D=p_S$ for some $L\ge2$.
Then $\spec_L(D)=\spec_L(S)$.
\end{corollary}

\begin{proof}[Proof sketch]
$p_D=p_S$ rewrites as $|S|\,\spec_L(D)=|D|\,\spec_L(S)$, so the two
integer spectrum vectors are proportional. By the lemma both have gcd
$1$ over their positive entries, and proportional gcd-one vectors are
equal (indeed $|S|\mid|D|$ and $|D|\mid|S|$, so $|S|=|D|$).
\end{proof}

\begin{remark}[Why primitivity alone is not enough]
The unrestricted statement --- primitivity plus $p_D=p_S$ implies
$\spec_L(D)=\spec_L(S)$ --- is false. At $L=2$, the primitive words
$S=AAB$ ($|S|=3$) and $D=AAABAB$ ($|D|=6$) have
$\spec_2(S)=(AA{:}1,AB{:}1,BA{:}1)$ and
$\spec_2(D)=(AA{:}2,AB{:}2,BA{:}2)$, hence identical normalized spectra
$p_S=p_D=(1/3,1/3,1/3)$ but distinct ordinary spectra. The lemma above
uses P2 essentially, via the BBT uniqueness input.
\end{remark}

\subsection{The complete-spectrum uniqueness input}

The remaining step is a classical one, supplied by an external source.

\begin{theorem}[Bresler--Bresler--Tse complete-spectrum uniqueness]
\label{thm:BBT}
\srcfact{} \cite{bresler2013}. Construct the $K$-mer graph from the complete
$(K+1)$-spectrum of a circular genome. If the genome satisfies Ukkonen's
condition --- no triple repeat and no interleaved repeat pair of length at least
$K$ --- then the graph has a unique Eulerian cycle, which spells the genome up
to cyclic rotation.
\end{theorem}

Setting $K=L-1$ aligns the theorem's condition exactly with P2: maximal triple
repeats then have length at most $L-2$, and every interleaved maximal-repeat pair
has a constituent of length at most $L-2$.

\subsection{Population uniqueness for the oriented primitive P2 class}

\begin{theorem}[Population uniqueness, oriented primitive P2]
\label{thm:population}
Let $L\ge2$. Among primitive P2-admissible oriented circular candidates, the
true genome $S$ is the unique population maximum-likelihood genome up to cyclic
rotation.
\end{theorem}

\begin{proof}
Let $D$ be primitive and P2-admissible with $\ellpop_S(D)=\ellpop_S(S)$. By
Lemma~\ref{lem:gibbs}, $p_D=p_S$; by Lemma~\ref{lem:scaling} both spectra
have gcd one, hence the corollary gives
$\spec_L(D)=\spec_L(S)$. Both $S$ and $D$ satisfy P2, hence Ukkonen's condition
at $K=L-1$, so Theorem~\ref{thm:BBT} applied to the complete $(L)$-spectrum
gives a unique Eulerian cycle and therefore $D\sim S$. The truth is a maximizer
by Lemma~\ref{lem:gibbs}.
\end{proof}

\statusline{The project-side reduction (division/Eulerian gcd-one, issue
\#70) and the thin final tie-to-rotation wiring (issue \#73) are
\kernelcheck{} in \texttt{PopulationReduction} and
\texttt{PopulationUniqueness}, reusing the \#70 reduction with no
duplicated infrastructure. The Gibbs/KL tie characterization and the BBT
uniqueness step remain explicit external premises (\srcfact{}
\cite{cover2006}, \cite{bresler2013}). The gcd-one/division writeup is the
project's.}

\subsection{Scope and one failed proof route}

The theorem is stated for the \emph{oriented} spectrum model. It must not be
silently transferred to reverse-complement-collapsed molecule classes, whose
observation object is different: population equality there is a statement about
class-indexed distributions, and the identification of classes changes which
words are distinguishable.

A stronger statement was attempted directly within the project, via a matching
argument on the transition involution of the length-$(L-1)$ de Bruijn structure.
That route has a genuine gap: a crossing between two raw occurrence pairs need
not survive maximal extension as two interleaved \emph{maximal} repeats, and an
explicit primitive example with the same length-$3$ spectrum and no cyclic
relation exhibits the failure (it nevertheless fails P2 for a different
interleaved pair, so it refutes the proof step, not the theorem). The
Bresler--Bresler--Tse route above bypasses this gap; the failed route is recorded
because it shows that ``the same spectrum is unique'' is not automatic from
short local repeat structure.

Finally, the population result does \emph{not} answer the finite 2016 question.
It answers a different, cleaner question: once empirical frequencies are replaced
by the true read distribution, bridging-style repeat bounds plus primitivity do
recover oriented circular uniqueness. Whether the historical finite-sample
statement holds under any fixed source-supported reading remains governed by
Section~\ref{sec:finite}.
